\Needspace{8\baselineskip}
\Article{201}
We add the following observations on the problem treated in the preceding articles.

1) We have shown how to solve \(tt-Duu=mm\) in every case where \(m\) is the greatest common divisor of three numbers \(M,2N,P\) such that \(NN-MP=D\). It is worth finding all numbers that can be such divisors, or all values of \(m\) for a given value of \(D\). Put \(D=nnD'\), with \(D'\) entirely free of square factors. We obtain this by taking \(nn\) to be the largest square dividing \(D\); if \(D\) itself already has no square factor, we must put \(n=1\). I then claim the following.

\textit{First}, if \(D'\) has the form \(4k+1\), every divisor of \(2n\) is a value of \(m\), and conversely. For if \(g\) divides \(2n\), we have the form \(\left(g,n,\frac{nn(1-D')}{g}\right)\), with determinant \(D\), in which \(g,2n,\frac{nn(D'-1)}{g}\) clearly have greatest common divisor \(g\) (for \(\frac{nn(D'-1)}{gg}=\frac{4nn}{gg}\cdot\frac{D'-1}{4}\) is clearly an integer). Conversely, suppose \(g\) is a value of \(m\), namely the greatest common divisor of \(M,2N,P\), with \(NN-MP=D\). Clearly \(4D\), or \(4nnD'\), is divisible by \(gg\). It follows that \(2n\) must be divisible by \(g\). For if \(g\) did not divide \(2n\), then \(g\) and \(2n\) would have a greatest common divisor smaller than \(g\). Putting this \(=\delta\), and
\SourceLeaf{198}
\(2n=\delta n',\ g=\delta g'\), we would have \(n'n'D'\) divisible by \(g'g'\), \(n'\) relatively prime to \(g'\), and therefore \(n'n'\) to \(g'g'\), so \(D'\) would also be divisible by \(g'g'\), contrary to the hypothesis that \(D'\) has no square factor.

\textit{Second}, if \(D'\) has the form \(4k+2\) or \(4k+3\), every divisor of \(n\) is a value of \(m\), and conversely every value of \(m\) divides \(n\). For if \(g\) divides \(n\), we have the form \(\left(g,0,-\frac{nnD'}{g}\right)\), whose determinant \(=D\), and in which \(g,0,\frac{nnD'}{g}\) clearly have greatest common divisor \(g\). --- Conversely, suppose \(g\) is a value of \(m\), namely the greatest common divisor of \(M,2N,P\), with \(NN-MP=D\). As above, \(g\) divides \(2n\), or \(\frac{2n}{g}\) is an integer. If this quotient were odd, the square \(\frac{4nn}{gg}\) would be \(\equiv1\ (\mathrm{mod.}\ 4)\), so \(\frac{4nnD'}{gg}\) would be either \(\equiv2\) or \(\equiv3\ (\mathrm{mod.}\ 4)\). But \(\frac{4nnD'}{gg}=\frac{4D}{gg}=\frac{4NN}{gg}-\frac{4MP}{gg}\equiv\frac{4NN}{gg}(\mathrm{mod.}\ 4)\), so \(\frac{4NN}{gg}\) would be either \(\equiv2\) or \(\equiv3\ (\mathrm{mod.}\ 4)\), \textit{a contradiction}, since every square must be congruent to zero or one modulo \(4\). So the quotient \(\frac{2n}{g}\) must be even, and \(\frac{n}{g}\) an integer, or \(g\) a divisor of \(n\).

Thus \(1\) is always a value of \(m\), or the equation \(tt-Duu=1\) can be solved by the preceding method for every positive nonsquare value of \(D\); and \(2\) is a value of \(m\) exactly when \(D\) has either the form \(4k\) or \(4k+1\).

2) If \(m\) is greater than \(2\) but is an admissible number, solving \(tt-Duu=mm\) can be reduced to solving a similar equation in which \(m\) is either \(1\) or \(2\). Put \(D=nnD'\) as before. If \(m\) divides \(n\), then \(mm\) divides \(D\). If the smallest values of \(p,q\) in \(pp-\frac{D}{mm}qq=1\) are \(p=P,\ q=Q\), the smallest values of \(t,u\) in \(tt-Duu=mm\) are \(t=mP,\ u=Q\). --- But if \(m\) does not divide \(n\), it at least divides \(2n\) and is certainly even, while \(\frac{4D}{mm}\) is an integer. If we then find that the smallest values of \(p,q\) in \(pp-\frac{4D}{mm}qq=4\) are \(p=P,\ q=Q\), the smallest values of \(t,u\) in \(tt-Duu=mm\) are \(t=\frac{m}{2}P,\ u=Q\). --- In both cases this method clearly derives not only from the smallest values of \(p,q\) the smallest values of \(t,u\), but from \textit{all} values of the former \textit{all} values of the latter.

3) Let \(t^0,u^0;\ t',u';\ t'',u''\) etc. denote all positive values of
\SourceLeaf{199}
\(t,u\) in \(tt-Duu=mm\) (as in the preceding article). Suppose some values in this sequence are congruent to its first values modulo any given \(r\), say \(t^\rho\equiv t\) (or \(\equiv m\)), \(u^\rho\equiv u^0\), or \(\equiv0\ (\mathrm{mod.}\ r)\), and at the same time the next values are congruent to the second values, say
\[
t^{\rho+1}\equiv t',\quad u^{\rho+1}\equiv u'\quad(\mathrm{mod.}\ r)
\]
Then we also have
\[
t^{\rho+2}\equiv t'',\quad u^{\rho+2}\equiv u'';\quad t^{\rho+3}\equiv t''',\quad u^{\rho+3}\equiv u'''\text{ etc.}
\]
This follows easily because both sequences \(t^0,t',t''\) etc., \(u^0,u',u''\) etc. satisfy recurrence relations: since
\[
t''=\frac{2T}{m}t'-t^0,\quad t^{\rho+2}=\frac{2T}{m}t^{\rho+1}-t^\rho
\]
we have
\[
t''\equiv t^{\rho+2}
\]
and similarly for the others. --- It follows that, in general,
\[
t^{h+\rho}\equiv t^h,\quad u^{h+\rho}\equiv u^h(\mathrm{mod.}\ r)
\]
where \(h\) denotes any number; and, more generally still, if
\[
\mu\equiv\nu\ (\mathrm{mod.}\ \rho),\quad\text{then }t^\mu\equiv t^\nu,\quad u^\mu\equiv u^\nu(\mathrm{mod.}\ r).
\]

4) We can always satisfy the conditions required in the preceding observation: we can always find an index \(\rho\) (for any given modulus \(r\)) such that
\[
t^\rho\equiv t^0,\quad t^{\rho+1}\equiv t',\quad u^\rho\equiv u^0,\quad u^{\rho+1}\equiv u'
\]
To prove this we observe the following.

\textit{First}, we can always satisfy the third condition. The criteria given in (1) readily show that \(pp-rrDqq=mm\) is also solvable. If the smallest positive values of \(p,q\) (apart from \(m,0\)) are \(P,Q\), the values of \(t,u\) clearly include \(t=P,\ u=rQ\). So \(P,rQ\) occur in the sequences \(t^0,t'\) etc., \(u^0,u'\) etc.; and if \(P=t^\lambda,\ rQ=u^\lambda\), then \(u^\lambda\equiv0\equiv u^0(\mathrm{mod.}\ r)\). It is also easy to see that between \(u^0\) and \(u^\lambda\) no term is congruent to \(u^0\) modulo \(r\).

\textit{Second}, if the other three conditions also hold, namely if \(u^{\lambda+1}\equiv u',\ t^\lambda\equiv t^0,\ t^{\lambda+1}\equiv t'\) as well, we clearly need only put \(\rho=\lambda\). If,
\SourceLeaf{200}
however, one or another of these conditions fails, I claim that we can certainly put \(\rho=2\lambda\). For equation~\([1]\) and the general formulas for \(t^e,u^e\) in the preceding article give
\[
t^{2\lambda}=\frac1m(t^\lambda t^\lambda+Du^\lambda u^\lambda)=\frac1m(mm+2Du^\lambda u^\lambda)
\]
and therefore
\[
\frac{t^{2\lambda}-t^0}{r}=\frac{2Du^\lambda u^\lambda}{mr}
\]
This quantity is an integer, since by hypothesis \(r\) divides \(u^\lambda\), and \(mm\) divides \(4D\), so all the more \(m\) divides \(2D\). --- Further, \(u^{2\lambda}=\frac2m t^\lambda u^\lambda\), and since
\[
4t^\lambda t^\lambda=4Du^\lambda u^\lambda+4mm
\]
which is therefore divisible by \(mm\), \(2t^\lambda\) is divisible by \(m\), so \(u^{2\lambda}\) is divisible by \(r\), or
\[
u^{2\lambda}\equiv u^0\ (\mathrm{mod.}\ r)
\]
Third, we find
\[
t^{2\lambda+1}=t'+\frac{2Du^\lambda u^{\lambda+1}}{m}
\]
and since \(\frac{2Du^\lambda}{mr}\) is an integer for the same reason, we have
\[
t^{2\lambda+1}\equiv t'\ (\mathrm{mod.}\ r)
\]
Finally, we find
\[
u^{2\lambda+1}=u'+\frac{2t^{\lambda+1}u^\lambda}{m}
\]
and since \(2t^{\lambda+1}\) is divisible by \(m\), and \(u^\lambda\) by \(r\), we have
\[
u^{2\lambda+1}\equiv u'\ (\mathrm{mod.}\ r).\quad \textit{This proves the claim.}
\]

We will see the use of the last two observations in what follows.

\Needspace{8\baselineskip}
\Article{202}
The special case of the problem, solving \(tt-Duu=1\), was already studied by the mathematicians of the preceding century. The very perceptive Fermat proposed this problem to the English analysts. Wallis names Brounker as the inventor of the solution he gives in \textit{Alg. Cap.}~98, \textit{Opp. T.}~II \textit{p.}~418 ff.; Ozanam names Fermat; and finally the distinguished Euler, who treats it in \textit{Comm. Petr.}~VI \textit{p.}~175, \textit{Comm. nov.}~XI \textit{p.}~28\HistoricalNote{*)}{In this paper the algorithm we presented in Article~27 is expressed in similar notation, which we neglected to mention there.}, \textit{Algebra P.}~II \textit{p.}~226, \textit{Opusc. An.}~I \textit{p.}~310, names Pell. This is why
\SourceLeaf{201}
some authors have called the problem \textit{Pell's problem}. All these solutions agree in substance with the one we obtain by choosing in Article~198 the reduced form in which \(a=1\). But no one before the distinguished La Grange rigorously\HistoricalNote{*)}{What Wallis offers for this purpose at the cited place, pp.~427, 428, carries no weight. His false reasoning consists in assuming, on p.~428, line~4, that for a given quantity \(p\) we can find integers \(a,z\) such that \(\frac{z}{a}\) is smaller than \(p\), with a shortfall smaller than the \textit{specified} one. This is certainly true when the \textit{specified} shortfall is \textit{a given quantity}, but not when it depends on \(a\) and \(z\), and so is variable, as happens in the present case.} proved that the operation they prescribe must eventually \textit{terminate}, or that the problem is always \textit{actually solvable}: \textit{Mélanges de la Soc. de Turin T.}~IV \textit{p.}~19, and more neatly in \textit{Hist. de l'Ac. de Berlin}, 1767, \textit{p.}~237. This investigation also appears in the \textit{supplements to Euler's Algebra}, already cited several times. Our method, however (based on entirely different principles and not restricted to \(m=1\)), usually offers several ways to reach the solution, since in Article~198 we can start from any other reduced form \((a,b,-a')\).
